% ID: SEL-D0055-OA-E04
% ESTADO: cotejada. Segunda pasada visual y borrador compilado revisados.
% Fuente: Selectividad/sel_2011_matematicas_aplicadas/122A-EXAMEN 6.pdf
% SHA-256: aebcbfddd4d5db5f5a10b4662e9e12b2ee362bbf574b751bccd6d78594766ad3
% Página: 1; opción A; ejercicio original 4.

El director de una televisión afirma que un nuevo programa que va a emitirse será visto, al menos, por un 30\% de personas. Una vez emitido se realizó una encuesta a 500 personas, elegidas al azar, y ésta reveló que 130 de ellas habían visto ese programa.
\begin{enumerate}[label=\alph*)]
\item \textbf{(0.5 puntos)} Formule la hipótesis nula y la alternativa del contraste de hipótesis que permite determinar si los datos de la encuesta realizada son compatibles con la afirmación del director.
\item \textbf{(1 punto)} Halle la región crítica de ese contraste para un nivel de significación del 5.5\%.
\item \textbf{(1 punto)} Según el dato obtenido en el apartado anterior \textquestiondown{}qué conclusión se obtiene sobre la afirmación realizada por el director de esa televisión?
\end{enumerate}
